\documentclass[11pt,a4paper]{article}
\usepackage[utf8]{inputenc}
\usepackage[T1]{fontenc}
\usepackage[spanish]{babel}
\usepackage{geometry}
\usepackage{amsmath}
\usepackage{amssymb}
\usepackage{graphicx}
\usepackage{enumitem}
\usepackage{booktabs}
\usepackage{array}
\usepackage{multirow}
\usepackage{fancyhdr}
\usepackage{xcolor}
\usepackage{mdframed}
\usepackage{tikz}
\usepackage{pgfplots}
\pgfplotsset{compat=1.18}

\geometry{top=2cm, bottom=2.5cm, left=2.5cm, right=2.5cm}

\pagestyle{fancy}
\fancyhf{}
\rhead{Tecnolog\'ia Digital VI --- 1er Semestre 2026}
\lhead{Parcial --- Modelo 2}
\rfoot{P\'agina \thepage}

\definecolor{tdblue}{RGB}{0,74,135}
\definecolor{lightgray}{RGB}{245,245,245}

\newmdenv[linecolor=tdblue, linewidth=1.5pt, roundcorner=4pt, backgroundcolor=lightgray, innertopmargin=6pt, innerbottommargin=6pt]{infobox}

\newcommand{\puntaje}[1]{\textbf{[#1 pts]}}

\begin{document}

%--- Encabezado ---
\begin{center}
  {\Large\bfseries Universidad Torcuato Di Tella}\\[4pt]
  {\large Tecnolog\'ia Digital VI --- Parcial Modelo 2}\\[2pt]
  {\normalsize 1er Semestre 2026}\\[8pt]
\end{center}

\begin{infobox}
\textbf{Instrucciones:} El parcial tiene una duraci\'on de \textbf{90 minutos}. Est\'a prohibido el uso de computadoras o c\'odigo. S\'i se permite el uso de una hoja de f\'ormulas. Cada secci\'on indica el puntaje. Puntaje total: \textbf{100 puntos}.
\end{infobox}

\vspace{0.5cm}

Nombre y apellido: \underline{\hspace{8cm}} \hfill Legajo: \underline{\hspace{3cm}}

\vspace{0.8cm}

%=============================================================
\section*{Secci\'on 1: Verdadero o Falso \puntaje{24}}
%=============================================================

\noindent Indique \textbf{V} (Verdadero) o \textbf{F} (Falso) para cada afirmaci\'on. Si la afirmaci\'on es falsa, \textbf{justifique brevemente} por qu\'e (sin justificaci\'on no se otorga puntaje en los falsos).

\vspace{0.3cm}

\begin{enumerate}[label=\textbf{\arabic*.}, leftmargin=*, itemsep=18pt]

\item \underline{\hspace{1cm}} En Random Forest, cada \'arbol se entrena con todas las features disponibles en cada split, lo cual garantiza la diversidad del ensamble.

\vspace{0.5cm}
\noindent\textit{Justificaci\'on (si es F):} \underline{\hspace{12cm}}

\vspace{0.2cm}
\underline{\hspace{14cm}}

\item \underline{\hspace{1cm}} En Gradient Boosting, aumentar el n\'umero de \'arboles ($M$) puede llevar a overfitting, a diferencia de Random Forest donde agregar m\'as \'arboles siempre es seguro.

\vspace{0.5cm}
\noindent\textit{Justificaci\'on (si es F):} \underline{\hspace{12cm}}

\vspace{0.2cm}
\underline{\hspace{14cm}}

\item \underline{\hspace{1cm}} Una de las ventajas de la regularizaci\'on L2 (Ridge) frente a L1 (Lasso) es que en presencia de variables correlacionadas, Ridge distribuye el peso entre ellas en lugar de seleccionar una sola.

\vspace{0.5cm}
\noindent\textit{Justificaci\'on (si es F):} \underline{\hspace{12cm}}

\vspace{0.2cm}
\underline{\hspace{14cm}}

\item \underline{\hspace{1cm}} El \textit{Partial Dependence Plot} (PDP) muestra el efecto marginal de un feature sobre la predicci\'on del modelo, promediado sobre todos los valores de los dem\'as features.

\vspace{0.5cm}
\noindent\textit{Justificaci\'on (si es F):} \underline{\hspace{12cm}}

\vspace{0.2cm}
\underline{\hspace{14cm}}

\item \underline{\hspace{1cm}} El bootstrap consiste en muestrear \textbf{sin} reemplazo del dataset original, garantizando que cada muestra aparezca en exactamente un subconjunto de entrenamiento.

\vspace{0.5cm}
\noindent\textit{Justificaci\'on (si es F):} \underline{\hspace{12cm}}

\vspace{0.2cm}
\underline{\hspace{14cm}}

\item \underline{\hspace{1cm}} PCA (An\'alisis de Componentes Principales) es un m\'etodo adecuado para capturar estructuras no lineales en los datos, como cl\'usteres en forma de espiral o c\'ırculo.

\vspace{0.5cm}
\noindent\textit{Justificaci\'on (si es F):} \underline{\hspace{12cm}}

\vspace{0.2cm}
\underline{\hspace{14cm}}

\item \underline{\hspace{1cm}} En AdaBoost, los clasificadores d\'ebiles que tienen un error ponderado mayor reciben un peso $\alpha_m$ m\'as grande en la predicci\'on final del ensamble.

\vspace{0.5cm}
\noindent\textit{Justificaci\'on (si es F):} \underline{\hspace{12cm}}

\vspace{0.2cm}
\underline{\hspace{14cm}}

\item \underline{\hspace{1cm}} El \textit{bias-variance tradeoff} establece que al aumentar la complejidad del modelo, el sesgo tiende a disminuir y la varianza tiende a aumentar, por lo que existe un punto \'optimo de complejidad que minimiza el error total.

\vspace{0.5cm}
\noindent\textit{Justificaci\'on (si es F):} \underline{\hspace{12cm}}

\vspace{0.2cm}
\underline{\hspace{14cm}}

\end{enumerate}

\vspace{0.5cm}

%=============================================================
\section*{Secci\'on 2: M\'ultiple Choice \puntaje{24}}
%=============================================================

\noindent Marque con un c\'irculo la \textbf{\'unica} opci\'on correcta en cada pregunta.

\vspace{0.3cm}

\begin{enumerate}[label=\textbf{\arabic*.}, leftmargin=*, itemsep=20pt]

\item ¿Cu\'al de las siguientes afirmaciones sobre XGBoost, LightGBM y CatBoost es \textbf{correcta}?

\begin{enumerate}[label=\alph*)]
  \item Los tres algoritmos son id\'enticos en su estrategia de crecimiento de \'arboles; solo difieren en el lenguaje de implementaci\'on.
  \item XGBoost crece el \'arbol nivel por nivel (\textit{level-wise}); LightGBM crece por hoja (\textit{leaf-wise}), siendo m\'as eficiente en memoria; CatBoost maneja variables categ\'oricas de forma nativa sin necesidad de encoding previo.
  \item LightGBM no puede manejar variables categ\'oricas y siempre requiere One-Hot Encoding.
  \item CatBoost utiliza expansi\'on de Taylor de segundo orden para optimizar su funci\'on objetivo, lo que lo hace m\'as preciso que XGBoost.
\end{enumerate}

\item ¿Qu\'e significa que dos componentes principales obtenidas por PCA sean ortogonales entre s\'i?

\begin{enumerate}[label=\alph*)]
  \item Que explican la misma cantidad de varianza en los datos.
  \item Que la correlaci\'on lineal entre ellas es cero; cada componente captura informaci\'on independiente de las dem\'as.
  \item Que sus autovalores asociados son iguales.
  \item Que no pueden ser usadas simult\'aneamente como predictores en un modelo supervisado.
\end{enumerate}

\item En el contexto del \textit{Bias-Variance Tradeoff}, ¿cu\'al de los siguientes modelos presenta \textbf{m\'as sesgo y menos varianza}?

\begin{enumerate}[label=\alph*)]
  \item Un \'arbol de decisi\'on con \texttt{max\_depth=20} y sin poda.
  \item Un Random Forest con 500 \'arboles.
  \item Una regresi\'on lineal aplicada a datos con relaciones no lineales complejas.
  \item Un modelo de Gradient Boosting con tasa de aprendizaje muy baja y muchos \'arboles.
\end{enumerate}

\item ¿Cu\'al es la raz\'on principal por la que Bagging reduce la varianza del modelo?

\begin{enumerate}[label=\alph*)]
  \item Porque penaliza los coeficientes del modelo con un t\'ermino de regularizaci\'on L2.
  \item Porque entrena los modelos secuencialmente, permitiendo que cada uno corrija los errores del anterior.
  \item Porque promedia las predicciones de m\'ultiples modelos entrenados sobre distintas muestras bootstrap, reduciendo la variabilidad que tendr\'ia cualquier modelo individual.
  \item Porque utiliza siempre el 100\% de los datos disponibles en cada modelo del ensamble.
\end{enumerate}

\item ¿Qu\'e ocurre con el error OOB (\textit{Out-of-Bag}) de un Random Forest a medida que se agregan m\'as \'arboles al ensamble?

\begin{enumerate}[label=\alph*)]
  \item Aumenta indefinidamente, indicando sobreajuste progresivo.
  \item Disminuye y se estabiliza, ya que agregar m\'as \'arboles a partir de cierto punto produce retornos decrecientes.
  \item Permanece constante independientemente de la cantidad de \'arboles.
  \item Primero disminuye y luego aumenta, al igual que la curva de validaci\'on en modelos de boosting.
\end{enumerate}

\item ¿Cu\'al de los siguientes m\'etodos de selecci\'on de features es un m\'etodo \textbf{Embedded}?

\begin{enumerate}[label=\alph*)]
  \item Correlaci\'on de Pearson entre cada feature y la variable objetivo.
  \item Forward Selection con validaci\'on cruzada.
  \item La importancia de features calculada por un \'arbol de decisi\'on o Random Forest durante el entrenamiento.
  \item Eliminar columnas con varianza menor a un umbral definido.
\end{enumerate}

\end{enumerate}

\vspace{0.5cm}

%=============================================================
\section*{Secci\'on 3: Interpretaci\'on de Gr\'aficos \puntaje{20}}
%=============================================================

\subsection*{Ejercicio A: Scree Plot de PCA \puntaje{10}}

El siguiente gr\'afico muestra el \textit{scree plot} de un PCA aplicado a un dataset con 8 features originales. Las barras representan la varianza explicada por cada componente principal (individual) y la l\'inea roja la varianza acumulada.

\vspace{0.3cm}
\begin{center}
\begin{tikzpicture}
\begin{axis}[
  width=10cm, height=6cm,
  xlabel={Componente Principal},
  ylabel={\% Varianza explicada},
  ymin=0, ymax=110,
  xmin=0.5, xmax=8.5,
  xtick={1,2,3,4,5,6,7,8},
  ytick={0,20,40,60,80,100},
  ybar,
  bar width=0.5cm,
  legend pos=north east,
  grid=major,
  grid style={line width=0.3pt, draw=gray!30},
  axis y line*=left,
]
\addplot[fill=tdblue!70, draw=tdblue] coordinates {
  (1,42)(2,24)(3,14)(4,8)(5,5)(6,4)(7,2)(8,1)
};
\addlegendentry{Varianza individual}
\end{axis}
\begin{axis}[
  width=10cm, height=6cm,
  ymin=0, ymax=110,
  xmin=0.5, xmax=8.5,
  xtick={1,2,3,4,5,6,7,8},
  axis y line*=right,
  axis x line=none,
  ylabel={\% Varianza acumulada},
]
\addplot[red, thick, mark=*, mark size=2pt] coordinates {
  (1,42)(2,66)(3,80)(4,88)(5,93)(6,97)(7,99)(8,100)
};
\addlegendentry{Var. acumulada}
\end{axis}
\end{tikzpicture}
\end{center}

\vspace{0.3cm}

\begin{enumerate}[label=\textbf{\alph*)}]
  \item Si se quisiera retener al menos el \textbf{85\%} de la varianza original, ¿cu\'antas componentes principales se deber\'ian seleccionar? Justifique leyendo el gr\'afico.

  \vspace{2.5cm}

  \item ¿Cu\'anto de la varianza original se \textbf{pierde} al reducir el dataset a 2 componentes principales? ¿Qu\'e implica esta p\'erdida en t\'erminos pr\'acticos?

  \vspace{2.5cm}

  \item El dataset original tiene 8 features y se quiere usar PCA antes de aplicar un \textbf{Random Forest}. Un compa\~nero sugiere que esto siempre mejora el modelo. ¿Est\'a de acuerdo? Justifique.

  \vspace{2.5cm}
\end{enumerate}

\subsection*{Ejercicio B: Partial Dependence Plot \puntaje{10}}

El siguiente gr\'afico muestra el PDP (\textit{Partial Dependence Plot}) del feature ``Edad del cliente'' sobre la predicci\'on de un modelo de churn (probabilidad de abandono).

\vspace{0.3cm}
\begin{center}
\begin{tikzpicture}
\begin{axis}[
  width=10cm, height=5.5cm,
  xlabel={Edad del cliente (a\~nos)},
  ylabel={Efecto marginal promedio (probabilidad de churn)},
  xmin=18, xmax=70,
  ymin=0.0, ymax=0.7,
  grid=major,
  grid style={line width=0.3pt, draw=gray!30},
]
\addplot[blue, thick, smooth] coordinates {
  (18,0.55)(22,0.50)(28,0.42)(35,0.32)(42,0.22)(50,0.18)(58,0.20)(65,0.25)(70,0.28)
};
\end{axis}
\end{tikzpicture}
\end{center}

\vspace{0.3cm}

\begin{enumerate}[label=\textbf{\alph*)}]
  \item Describa en palabras qu\'e muestra este gr\'afico. ¿C\'omo var\'ia la probabilidad de churn en funci\'on de la edad?

  \vspace{2.5cm}

  \item ¿El PDP permite concluir que la edad \textbf{causa} el churn? ¿Por qu\'e?

  \vspace{2.5cm}

  \item Si en lugar del PDP se mostrara un \textbf{ICE plot} (\textit{Individual Conditional Expectation}), ¿qu\'e informaci\'on adicional aportar\'ia que el PDP no muestra?

  \vspace{2.5cm}
\end{enumerate}

\vspace{0.3cm}

%=============================================================
\section*{Secci\'on 4: Desarrollo \puntaje{32}}
%=============================================================

\begin{enumerate}[label=\textbf{\arabic*.}, leftmargin=*, itemsep=10pt]

\item \puntaje{10} \textbf{Ensambles y Diversidad}

\begin{enumerate}[label=\textbf{\alph*)}]
  \item Explique por qu\'e la diversidad entre los modelos base es fundamental para que un ensamble mejore el rendimiento respecto a un modelo individual. ¿Qu\'e suceder\'ia si todos los modelos del ensamble cometieran exactamente los mismos errores?
  \vspace{3.5cm}
  \item Random Forest introduce dos fuentes de aleatoriedad para lograr diversidad. Identif\'iquelas y explique c\'omo cada una contribuye a reducir la correlaci\'on entre \'arboles.
  \vspace{3.5cm}
  \item ¿Por qu\'e Bagging es trivialmente paralelizable pero Boosting no lo es? Explique la diferencia estructural entre ambos m\'etodos que determina esto.
  \vspace{3cm}
\end{enumerate}

\item \puntaje{10} \textbf{Reducci\'on de la Dimensionalidad}

\begin{enumerate}[label=\textbf{\alph*)}]
  \item Explique el concepto de ``maldici\'on de la dimensionalidad''. ¿Por qu\'e los algoritmos basados en distancias (como KNN) se ven especialmente perjudicados en espacios de alta dimensi\'on?
  \vspace{3.5cm}
  \item Compare PCA con t-SNE y UMAP en t\'erminos de: (i) tipo de relaciones que capturan (lineales vs. no lineales), (ii) capacidad de proyectar nuevas muestras, y (iii) utilidad para preprocesamiento vs. solo visualizaci\'on.
  \vspace{4cm}
  \item ¿Por qu\'e es obligatorio estandarizar los features \textbf{antes} de aplicar PCA? ¿Qu\'e problema ocurrir\'ia si no se hace?
  \vspace{3cm}
\end{enumerate}

\item \puntaje{12} \textbf{Caso integral: Detecci\'on de piezas defectuosas}

Una f\'abrica de electr\'onicos quiere predecir qu\'e piezas del proceso de manufactura saldr\'an defectuosas antes de que lleguen al control de calidad. Cuentan con:
\begin{itemize}
  \item Datos de sensores durante la fabricaci\'on: temperatura, presi\'on, vibraci\'on (200 sensores en total).
  \item Registros hist\'oricos de 500.000 piezas, de las cuales 4.500 resultaron defectuosas.
  \item Algunos sensores tienen datos correlacionados entre s\'i.
\end{itemize}

\begin{enumerate}[label=\textbf{\alph*)}]
  \item ¿Qu\'e desaf\'ios presenta este dataset y c\'omo los abordar\'ia? Considere el desbalance de clases, la alta dimensi\'on y la correlaci\'on entre sensores.
  \vspace{4cm}
  \item ¿Qu\'e modelo(s) elegiria para este problema y por qu\'e? Compare al menos dos opciones considerando: rendimiento predictivo, interpretabilidad y complejidad de ajuste.
  \vspace{3.5cm}
  \item En una f\'abrica, los errores Falso Negativo (no detectar una pieza defectuosa) y Falso Positivo (descartar una pieza buena) tienen costos muy distintos. ¿C\'omo tendr\'ia esto en cuenta al elegir el umbral de decisi\'on y la m\'etrica de evaluaci\'on?
  \vspace{3.5cm}
\end{enumerate}

\end{enumerate}

\vspace{1cm}
\begin{center}
\textit{--- Fin del Parcial Modelo 2 ---}
\end{center}

\end{document}
